\chapter{Introdução} \label{cap:intro}

 O agendamento eficiente de recursos para a execução de tarefas cumpre um papel fundamental em ambientes de produção e manufatura, bem como em indústrias de transporte e distribuição. Nesse sentido, diferentes cenários podem exigir a otimização de diferentes objetivos. A partir disso, podemos modelar tais cenários de agendamento de tarefas como problemas de otimização combinatória. A flexibilidade desses problemas nos permite acoplar diversas características à modelagem, de forma que ela represente fielmente um cenário real \citep{pinedo2008}. A partir dessa abordagem, podemos representar uma infinidade de processos de gerenciamento de recursos.

% A formalização de eficiência e a criação de ferramentas para classificar problemas pela sua dificuldade inerente estão entre os avanços mais importantes da computação. \cite{eva2005} definem como eficientes, algoritmos que executam em um tempo limitado por um polinômio no tamanho da entrada recebida. A partir dessa ideia, podemos caracterizar o conjunto de problemas para os quais conhecemos algoritmos exatos eficientes $P$. Dentro do universo de problemas fora da classe $P$, uma classe que se destaca pela sua importância são os problemas NP-Difíceis. A comunidade científica acredita que não há maneira eficiente de resolver os problemas dessa classe, pois, pela sua definição, se encontrássemos uma solução exata eficiente para algum deles, todos os outros problemas da classe poderiam ser resolvidos eficientemente.

 Um caso específico de problema de agendamento é o problema \ac{JITJSS}. Além da alocação de máquinas para executar tarefas, esse problema define prazos para cada etapa de uma tarefa e caso alguma dessas etapas seja finalizada fora do seu prazo de entrega, ela gera uma penalidade de adiantamento ou de atraso. O objetivo ao resolver esse problema é encontrar uma alocação que minimize essas penalidades. Esse problema foi proposto pela primeira vez por \cite{baptiste2008lagrangian}, que também propuseram as instâncias do problema, juntamente com os primeiros limites superiores e inferiores para essas instâncias.

A relevância do \ac{JITJSS} reside na sua capacidade de modelar cenários nos quais a finalização de processos dentro de um prazo é essencial. Por exemplo, em uma malha ferroviária, alocar os tempos em que cada locomotiva ocupa cada estação é essencial para o fluxo eficiente dos trens. Podemos modelar as paradas como tarefas que precisam de recursos (estações) para serem executadas e que duram um determinado tempo. Além disso, cada parada de trem possui um prazo desejável para encerrar e liberar a estação. Dessa maneira, queremos alocar as estações minimizando ao máximo os desvios desses prazos de liberação. Para mais detalhes sobre o problema de alocação de trens, veja \cite{liu2011scheduling} e \cite{flamini2008real}.

Considere, como outro exemplo de aplicação, os procedimentos cirúrgicos de um hospital. Para a execução desses procedimentos, há recursos limitados disponíveis que serão alocados por um determinado tempo para suas diferentes fases. Os processos de cada paciente possuem tempos de duração diferentes e uma ordem de execução predefinida. Também existem prazos desejáveis para que cada processo seja finalizado. Dessa forma, uma maneira de alocar tais processos hospitalares aos recursos disponíveis, minimizando as execuções fora do prazo, otimiza os processos hospitalares descritos. Veja \cite{pham2008surgical} para mais detalhes sobre esse problema. Embora essas aplicações do mundo real possuam suas próprias complexidades, elas ilustram a estrutura e a importância dos problemas de agendamento com restrições de tempo, como o \ac{JITJSS}.

Além de representar características de cenários do mundo real, o problema \ac{JITJSS} é fortemente NP-Difícil \citep{du1990minimizing}. Para problemas dessa classe, não existem soluções exatas eficientes, a menos que P$=$NP \citep{cormen2022introduction}. Em vista disso, uma opção é tentar encontrar soluções exatas, embora ineficientes. O problema nessa abordagem é que encontrar soluções exatas se torna computacionalmente inviável para instâncias grandes. Uma outra opção é trocar a otimalidade das soluções por um algoritmo mais eficiente. Nesse sentido, existem algoritmos aproximados que encontram soluções garantidamente próximas do ótimo. Por exemplo, um algoritmo 1,2-aproximado para um problema de minimização sempre encontra soluções até 20\% piores do que a solução ótima. O problema nessa abordagem é que desenvolver tais estratégias exige um entendimento pleno das características do problema e demonstrações robustas da aproximação construída. Por esses motivos, uma opção recorrente para resolver problemas difíceis é utilizar heurísticas \citep{eva2005}.

Heurísticas são soluções que não garantem a otimalidade da solução em troca de obter soluções em tempo viável. Por não terem o compromisso com a otimalidade, desenvolver estratégias heurísticas livra o projetista dos problemas no desenvolvimento de métodos exatos e aproximados \citep{eva2005}.

Podemos encontrar na literatura diversas aplicações de heurísticas ao problema \ac{JITJSS}. A título de exemplo, temos algoritmos evolucionários \citep{araujo2009genetic,dos2010combination, yang2012improved}, a busca em grandes vizinhanças \citep{laborie2007self, abderrazzak2022adaptive} e a busca em vizinhança variável \citep{wang2014variable, ahmadian2021meta}.

A eficácia de uma heurística, no entanto, depende de uma série de decisões de projeto que os projetistas precisam realizar. Essas decisões são baseadas em ciclos de desenvolvimento e validação das estratégias desenvolvidas. Porém, mesmo com o processo de desenvolvimento finalizado, algumas escolhas de implementação permanecem em aberto, e o projetista pode expor essas decisões aos usuários através de parâmetros do algoritmo. Dessa forma, o usuário pode definir a combinação de parâmetros que obtém os melhores resultados para o seu objetivo. O problema é que encontrar tal combinação de parâmetros pode se tornar uma tarefa extremamente custosa à medida que o número de parâmetros cresce \citep{birattari2009tuning}.

Nesse sentido, definimos uma configuração de um algoritmo como uma atribuição de valores a cada parâmetro. A partir disso, temos que o processo de realizar a busca da configuração ideal manualmente possui alguns problemas. Além de ser um processo custoso e tedioso para o indivíduo responsável, é um processo muito propenso a erros, dificilmente reproduzível e, mesmo com muito esforço do configurador, ainda é um processo inconclusivo pois não nos dá argumentos suficientes de que a configuração encontrada é realmente boa. Para resolver esse problema, a comunidade científica formalizou a configuração de algoritmos e desenvolveu diversos métodos de Configuração Automática de Algoritmos (do inglês, \textit{Automatic Algorithm Configuration} -- \acs{AAC}) \citep{birattari2009tuning}. Esses métodos nos permitem explorar o espaço de configurações sistematicamente e, com isso, encontrar configurações comprovadamente boas.

Utilizando as técnicas de \ac{AAC}, fica mais fácil encontrar boas configurações para um algoritmo. Ainda assim, primeiro precisamos construir um algoritmo para conseguirmos configurá-lo. Para isso, precisamos testar diversas combinações de componentes algorítmicos para encontrarmos bons algoritmos que resolvam um certo problema. Esse processo, se realizado de forma manual, sofre dos mesmos problemas da configuração manual de algoritmos. Levando a ideia do \ac{AAC} um passo adiante e com o objetivo de otimizar esse processo, a comunidade científica propôs expor totalmente as decisões de projeto de algoritmos às técnicas de \ac{AAC}. Dessa maneira, seria possível explorar sistematicamente o espaço de algoritmos para um problema. A partir dessa ideia, a comunidade científica desenvolveu diversos métodos e conceitos de Projeto Automático de Algoritmos (do inglês, \textit{Automatic Algorithm Design} -- \acs{AAD}) \citep{stutzle2018automated}.

Com o avanço das técnicas de \ac{AAD}, esses métodos vêm sendo utilizados para gerar, automaticamente, novos algoritmos estado da arte para diferentes problemas de otimização combinatória, como o Problema da Satisfabilidade Proposicional (do inglês, \textit{Boolean Satisfiability Problem -- \acs{SAT}}) \citep{khudabukhsh2016satenstein}, \ac{UBQP} \citep{de2018automatic, de2018automatically}, \ac{FSP} \citep{brum2018automatic}, \ac{PFSP} \citep{marmion2013automatic, pagnozzi2019automatic}, \ac{HFS} \citep{alfaro2020automatic} e \ac{NPFSP} \citep{brum2022automatic}.

A partir de todas essas ideias, percebemos o sucesso na aplicação do \ac{AAD} em diversos problemas de otimização, incluindo vários problemas de agendamento. Contudo, ainda não há na literatura um trabalho que aplique as técnicas de \ac{AAD} para resolver o problema \ac{JITJSS}. Por isso, considerando a relevância do \ac{JITJSS} e a lacuna na literatura, o presente trabalho parte da hipótese de que a aplicação das técnicas de projeto automático tenha a capacidade de gerar novos algoritmos que superem o estado da arte atual. 

\section{Objetivos}

\subsection{Objetivo Geral}

O presente trabalho tem como objetivo aplicar as técnicas de \ac{AAC} e \ac{AAD} na resolução do problema \ac{JITJSS} e, com isso, contribuir para um melhor entendimento sobre o problema e sobre a aplicação das técnicas de \ac{AAC} e \ac{AAD} em problemas de agendamento.

\subsection{Objetivos Específicos}

Para atingir o nosso objetivo geral, estabelecemos os seguintes objetivos específicos:

\begin{enumerate}
  \item Definir o espaço de projeto de algoritmos.
  \item Propor pelo menos um componente heurístico específico para o \ac{JITJSS}, a fim de incrementar o espaço de projeto de algoritmos.
  \item Aplicar as técnicas de \ac{AAC} e \ac{AAD} sobre o espaço de projeto definido anteriormente, adotando as boas práticas metodológicas da literatura para garantir a robustez e a validade do processo de configuração.
  \item Avaliar a qualidade das estratégias obtidas pela aplicação das técnicas de \ac{AAC} e \ac{AAD}.
  % \item Fornecer uma comparação entre os algoritmos gerados automaticamente e os algoritmos estado da arte para o \ac{JITJSS}.
  % \item Validar a efetividade dos componentes da literatura em resolver o problema por meio dos resultados da aplicação dos métodos de \ac{AAC} e \ac{AAD}.
\end{enumerate}

\section{Organização do Trabalho}

Além deste capítulo introdutório, o presente trabalho está organizado da seguinte forma: O Capítulo \ref{chapt:ref_teorico} detalha os conceitos que embasam o presente trabalho; O Capítulo \ref{chapt:trabalhos} realiza uma análise da literatura relacionada ao presente trabalho; O Capítulo \ref{chapt:metodologia} explica os métodos e materiais que serão utilizados na realização deste trabalho; O Capítulo \ref{chapt:resultados} expõe o trabalho realizado e os resultados obtidos; E, por fim, o Capítulo \ref{chapt:conclusao} consolida os resultados obtidos e apresenta os possíveis trabalhos futuros.


